\section{Introduction}
\label{sec:introduction}

Offline reinforcement learning (RL) optimizes a policy from a fixed dataset and therefore cannot repair value errors by collecting corrective experience. In deterministic actor--critic methods, policy improvement and critic learning are tightly coupled: the actor is trained to increase a learned value function, and a target copy of that actor supplies the next action in the critic's Bellman target. Increasing actor aggressiveness can improve a conservative policy, but it can also move the target action into regions where the critic is extrapolating. The resulting error is then recycled into subsequent actor and critic updates \citep{fujimoto2019offpolicy,kumar2019bear,kumar2020cql}.

TD3+BC controls this feedback by combining value maximization with a squared behavior-cloning penalty \citep{fujimoto2021minimalist}. Its scalar actor coefficient nevertheless controls two conceptually different quantities: the displacement of the branch whose target copy enters bootstrapping and the reach of the policy used at deployment. A one-branch architecture cannot vary these quantities separately.

We study \emph{multi-step proximal policy improvement} (MPI), a persistent chain of $K$ deterministic actors. A total nominal budget $T$ is split into local coefficients $h=T/K$. The first actor is anchored to dataset actions, later actors are anchored to stop-gradient outputs of their predecessors, a target copy of the first actor enters the TD target, and the final actor is evaluated. Thus, increasing $K$ simultaneously makes the critic-coupled update more local and extends deployment through additional actor branches.

This construction has a useful geometric interpretation. For Dirac policies, squared action displacement is exactly a quadratic Wasserstein cost. The corresponding idealized actor objective is a proximal step for the negative learned critic. This interpretation yields two limited but precise statements: an approximately solved proximal subproblem satisfies a descent inequality relative to the learned critic, and the one-step perturbation of a Lipschitz TD target depends on the first-hop displacement rather than the final-hop displacement. Neither statement proves monotone environment return. The released implementation also performs one Adam update per actor instead of solving each proximal subproblem, so we refer to it as a \emph{proximal-loss realization} rather than an exact implicit integrator.

The empirical object is a fixed-budget phase diagram, not a conventional single-hyperparameter benchmark. We sweep fourteen budgets over nine D4RL locomotion tasks. The original two-seed grid contains TD3+BC ($K=1$), proximal MPI for $K\in\{2,3\}$, and scale-matched linearized MPI for $K\in\{2,3\}$. A newly completed proximal $K=4$ grid extends the comparison. Three observations are robust in the available matrices. First, additional proximal actors primarily improve the lower tail: the high-budget mean rises from $25.97$ at $K=1$ to $63.61$ at $K=4$, while collapsed environment--budget cells fall from $35/63$ to $9/63$. Second, $K=4$ improves over $K=3$ by $10.83$ points on average but only $0.32$ in the median, again indicating failure rescue rather than a uniform shift. Third, a complete linearized $K=3$ sweep improves over linearized $K=2$, showing that the empirical effect is not unique to the proximal loss.

The expanded artifact also changes the appropriate scientific claim. Increasing $K$ changes local actor coefficients, actor compute, target/deployment routing, persistent network initialization, and the critic trajectory. Moreover, the current code generates the critic key as the last element of a $K$-dependent random-key split, so nominally paired seeds do not share critic initialization across $K$. We therefore interpret the sweep as evidence for an end-to-end training bundle and identify a $K$-invariant initialization rerun as necessary before assigning a causal effect to budget splitting itself.

Our contributions are:
\begin{itemize}
    \item We formulate MPI as a target/deployment actor chain and give a deterministic fiberwise Wasserstein interpretation of its squared anchors.
    \item We state the theoretical guarantees at their actual level: conditional learned-critic descent and first-hop target perturbation, not true-return monotonicity or an exact JKO solve.
    \item We extend the complete fixed-budget study to proximal $K=4$, where the high-budget mean reaches $63.61$ and the collapse count falls to $9/63$.
    \item We audit the released implementation and separate descriptive sweep evidence from unresolved causal and reproducibility limitations, most notably $K$-dependent critic initialization and missing raw diagnostic logs.
\end{itemize}
